\documentclass[10pt]{article}
\usepackage[T1]{fontenc}
\usepackage[margin=0.78in]{geometry}
\usepackage{amsmath,amssymb,amsthm}
\usepackage{graphicx}
\usepackage[colorlinks=true,linkcolor=blue,citecolor=blue,urlcolor=blue]{hyperref}
\newtheorem{theorem}{Theorem}
\newtheorem{lemma}{Lemma}
\newcommand{\B}{\mathbf B}
\newcommand{\Rcal}{\mathcal R}
\newcommand{\E}{\mathcal E}
\newcommand{\K}{\mathcal K}
\newcommand{\N}{\mathbb N}
\title{A positive Poisson-representation theorem\\with one nonscalar row}
\author{Run \texttt{fa\_banach\_001} --- \texttt{agent\_lane\_11}}
\date{13 August 2026}
\begin{document}
\maketitle

\noindent\textbf{Status.} Verified partial result; likely valid subject to
human review.  It gives a new dimension class in which the open question of
Popescu~\cite{source} has an affirmative answer.

\section{The question and the result}

For the regular polyball $\B_{\mathbf n}$, $\mathbf n=(n_1,\ldots,n_k)$,
Popescu asks whether every positive free $k$-pluriharmonic function is the
Poisson transform of a completely positive map on the $C^*$-algebra generated
by the right universal model.

\begin{center}
\includegraphics[width=.94\textwidth]{figures/source_question_crop.png}
\end{center}

The cited Proposition 3.8 actually proves $k=1$ and the polydisk case
$n_1=\cdots=n_k=1$.  (The last displayed source sentence appears to omit the
final ``$=1$''.)  The following strictly enlarges those proved cases.

\begin{theorem}\label{thm:main}
Let $\mathbf n\in\N^k$ have at most one coordinate larger than one.  Every
positive free $k$-pluriharmonic function
\[
 F:\B_{\mathbf n}(\mathcal H)\longrightarrow
 B(\E)\otimes_{\min}B(\mathcal H)
\]
is the noncommutative Poisson transform $\boldsymbol{\mathcal P}\mu$ of a
completely positive map $\mu:C^*(\mathbf R)\to B(\E)$.
\end{theorem}

Thus arbitrary operator-valued coefficients and arbitrarily many scalar
coordinates are allowed, together with one genuinely free row of any finite
size.

\section{The simultaneous unitary-extension lemma}

The only new dilation input is elementary but useful.

\begin{lemma}\label{lem:dilation}
Let $U_1,\ldots,U_s$ be commuting isometries on $\K$, and let
$V=(V_1,\ldots,V_n)$ be a row isometry on $\K$ such that every $V_j$
commutes with every $U_i$.  There are a Hilbert space
$\widetilde\K\supset\K$, commuting unitaries
$\widetilde U_1,\ldots,\widetilde U_s$, and a row isometry
$\widetilde V=(\widetilde V_1,\ldots,\widetilde V_n)$ extending the original
operators, such that the tuple
$(\widetilde V,\widetilde U_1,\ldots,\widetilde U_s)$ is doubly commuting.
\end{lemma}

\begin{proof}
For $p\le q$ in $\N^s$, use the isometry
$U^{q-p}=U_1^{q_1-p_1}\cdots U_s^{q_s-p_s}$ as the connecting map from a
copy $\K_p$ of $\K$ to $\K_q$.  Let $\widetilde\K$ be the Hilbert direct
limit and $J_p:\K\to\widetilde\K$ its canonical isometries, normalized by
\[
 J_p=J_qU^{q-p}\qquad(p\le q).
\]
Set, initially on the dense union of the ranges of the $J_p$,
\[
 \widetilde U_iJ_ph=J_pU_ih,
 \qquad
 \widetilde V_jJ_ph=J_pV_jh.
\]
The definitions are independent of the representative because all original
operators commute.  Inner products show that the $\widetilde V_j$ are
isometries with mutually orthogonal ranges.  Each $\widetilde U_i$ is an
isometry, and it is onto because
$\widetilde U_iJ_{p+e_i}h=J_ph$; hence it is unitary.  The extended operators
commute.  Fuglede's theorem now gives
$\widetilde V_j\widetilde U_i^*=\widetilde U_i^*\widetilde V_j$.  The
commuting unitaries also commute with one another's adjoints, so the rows are
doubly commuting.
\end{proof}

\section{Proof of Theorem \ref{thm:main}}

First suppose $F(0)=I_{\E}$.  The positive-structure theorem in Section 3 of
\cite{source} supplies a Hilbert space $\K\supset\E$ and commuting row
isometries $\mathbf V=(V_1,\ldots,V_k)$ whose compressed moments are exactly
the coefficients of $F$:
\[
 A_{(\boldsymbol\alpha;\boldsymbol\beta)}
 =P_{\E}\mathbf V_{\widetilde{\boldsymbol\alpha}}^*
       \mathbf V_{\widetilde{\boldsymbol\beta}}\!mid_{\E}.
\tag{1}
\]
After relabeling, $V_1$ is the possible nonscalar row, while
$V_2,\ldots,V_k$ are single commuting isometries.  Lemma~\ref{lem:dilation}
extends this tuple to a doubly commuting row-isometric tuple
$\widetilde{\mathbf V}$ on $\widetilde\K$.

Because $\widetilde{\mathbf V}$ extends $\mathbf V$, its analytic products
agree with those of $\mathbf V$ on $\K$.  Hence all compressed moments in
(1) are unchanged.  Double commutation makes
\[
 \pi(\mathbf R_{\boldsymbol\alpha}
              \mathbf R_{\boldsymbol\beta}^*)
 :=\widetilde{\mathbf V}_{\boldsymbol\alpha}
   \widetilde{\mathbf V}_{\boldsymbol\beta}^*
\]
a $*$-representation of $C^*(\mathbf R)$.  If
$J:\E\hookrightarrow\widetilde\K$ is inclusion, then
$\mu(T)=J^*\pi(T)J$ is completely positive.  Its moments are (1), so
uniqueness of the free pluriharmonic coefficients gives
$F=\boldsymbol{\mathcal P}\mu$.

For general $A:=F(0)\ge0$ and $\varepsilon>0$, normalize
\[
 G_\varepsilon=(A+\varepsilon I)^{-1/2}
 (F+\varepsilon I)(A+\varepsilon I)^{-1/2}.
\]
The normalized case, followed by conjugation with
$(A+\varepsilon I)^{1/2}$, produces a completely positive map
$\mu_\varepsilon$ representing $F+\varepsilon I$ and satisfying
$\|\mu_\varepsilon\|=\|A+\varepsilon I\|$.  A point-ultraweak cluster limit
as $\varepsilon\downarrow0$ is completely positive.  Its moments converge to
the coefficients of $F$, and therefore it represents $F$.  This proves the
theorem. \qed

\section{Why the full arbitrary-row case remains}

The proof uses that the semigroup generated by the scalar rows is $\N^s$:
its isometric representation has the direct-limit unitary extension above.
With two nonscalar rows one would need to Cuntz-dilate one free row while
extending the other commuting row and preserving its orthogonal-range
relations.  There is no analogous Ore direct limit, and the defect space of a
row isometry need not be invariant under the other row.  General
product-system regular-dilation theorems impose additional positivity or
regularity hypotheses.  Consequently the argument does not settle the
original question for two or more genuinely free rows.

A bounded search using the exact question and this dilation formulation found
no later primary-source answer or this intermediate dimension theorem.  The
novelty claim is therefore only \emph{apparently new within the bounded
search}.  Eight focused upgrade attempts, including finite polynomial
counterexample searches and an attempted iterative Cuntz dilation, are
recorded in the accompanying attempt file.

\begin{thebibliography}{1}
\bibitem{source}
G.~Popescu, \emph{Free pluriharmonic functions on noncommutative polyballs},
Analysis \& PDE \textbf{9} (2016), 1185--1234;
arXiv:1511.08852.
\end{thebibliography}
\end{document}

